% =============================================================================
% notas_reunion_hoy.tex — Notas personales de una página para la reunión Zoom
% con Gonzalo. NO es un guion para presentarle a él; es para tener al lado.
% Espíritu: abrir, mostrar, preguntar, callar. Seguirlo si se desvía.
% Compilar: pdflatex notas_reunion_hoy.tex
% =============================================================================
\documentclass[11pt,a4paper]{article}

\usepackage[utf8]{inputenc}
\usepackage[T1]{fontenc}
\usepackage[spanish,es-nodecimaldot]{babel}
\usepackage[margin=2cm]{geometry}
\usepackage{enumitem}
\usepackage{xcolor}
\usepackage{titlesec}

\titleformat{\section}{\normalfont\bfseries\large}{}{0pt}{}
\titlespacing{\section}{0pt}{8pt}{4pt}
\setlength{\parskip}{2pt}
\pagestyle{empty}

\begin{document}

{\Large\bfseries Reunión con Gonzalo — hoy}\\[2pt]
{\small Un solo objetivo: que él haga el análisis conmigo, no que yo lo presente cerrado.}

\section*{Regla de la reunión}
No uso las 15 láminas como guion. Las abro \emph{solo si él pide} profundizar
en algo puntual. Si se queda media hora en un punto, lo sigo ahí —
no vuelvo a mi orden.

\section*{Orden}

\begin{enumerate}[leftmargin=1.4em,itemsep=4pt]
  \item \textbf{Abrir con la imagen} que ya le gustó (secuencia de 4 paneles:
  caída, compresión, expansión, apogeo). Sin conclusiones todavía —
  solo ``mira lo que se ve''.

  \item \textbf{Reproducción 2012: una frase.} ``Reproducimos su trabajo,
  coincide en el régimen elástico.'' Punto. No me quedo ahí.

  \item \textbf{Mostrar el dato crudo y preguntar, no anunciar.}\\
  \emph{``Medimos $N_c$ de tres formas distintas y dan 25, 60 y 350 —
  ¿cómo lo lees tú?''}\\
  Dejar que él arme la historia de las tres escalas antes de decir
  la mía.

  \item \textbf{Silencio después de preguntar.} Contar hasta cinco antes
  de completar la idea.

  \item \textbf{Las 4 decisiones (autoría, ORCID, idioma, arXiv/RBEF)
  van al final, solo si se llega ahí solo.} No las fuerzo hoy.
\end{enumerate}

\section*{Si me pregunta ``¿y qué sigue?''}
Ahí sí puedo mostrar la reformulación ($f_{ball}\to 1$, $N_c^{ef}\approx25$)
como algo abierto a su lectura, no como resultado cerrado.

\section*{Recordatorio}
La vez pasada avancé más de lo pedido y él, con calma, me devolvió al
problema del rebote. Hoy: seguirlo a él, no a la agenda escrita.

\end{document}
